%!Mode:: "TeX:UTF-8"
%!TEX encoding = UTF-8 Unicode
%lan
%arara: xelatex
\documentclass{ctexart}
% \usepackage{listings}

\newif\ifpreface
%\prefacetrue
\input{../../../global/all}
\begin{document}
\large
\setlength{\baselineskip}{1.2em}
\ifpreface
\input{../../../global/preface}
\newgeometry{left=2cm,right=2cm,top=2cm,bottom=2cm}
\else
\newgeometry{left=2cm,right=2cm,top=2cm,bottom=2cm}
\maketitle
\fi
%from_here_to_type

\section{Problem 20.18}

Let
\[
  A=
  \begin{bmatrix}
    0.1000 & 0.5000 & -0.1000\\
    0.4000 & 0.2000 & 0.6000\\
    0.2000 & -0.3000 & 0.4000
  \end{bmatrix}.
\]
Since $A$ is not diagonally dominant, the convergence of the Jacobi and Gauss--Seidel methods is not guaranteed.

\paragraph{(a) Norms of $B_J$ and $B_{GS}$}

The computed norms of $B_J$ are
\[
  \|B_J\|_1 = 5.7500,\qquad
  \|B_J\|_2 = 5.2180,\qquad
  \|B_J\|_\infty = 6.0000.
\]

The computed norms of $B_{GS}$ are
\[
  \|B_{GS}\|_1 = 25.0000,\qquad
  \|B_{GS}\|_2 = 16.3566,\qquad
  \|B_{GS}\|_\infty = 15.0000.
\]

\paragraph{(b) Spectral radii}

\[
  \rho(B_J)=3.1756,\qquad
  \rho(B_{GS})=3.7500.
\]

\paragraph{(c) Convergence analysis}

A stationary iteration converges if and only if
\[
  \rho(B)<1.
\]
Since
\[
  \rho(B_J)>1,\qquad \rho(B_{GS})>1,
\]
neither the Jacobi method nor the Gauss--Seidel method converges.

To verify this numerically, we choose a random vector $b$ and set the initial guess
\[
  x^{(0)}=0.
\]
The tolerance is chosen as
\[
  \text{tol} = 10^{-6},
\]
and the maximum number of iterations is set to
\[
  \text{numiter} = 2{,}000{,}000.
\]

Despite allowing a large number of iterations and a reasonable tolerance, both methods fail to converge:
\[
  \text{GS does not converge},\qquad
  \text{Jacobi does not converge}.
\]

This agrees with the theoretical prediction, since the spectral radii of both iteration matrices are greater than $1$.

\paragraph{Conclusion}

Both the Jacobi and Gauss--Seidel methods do not converge for this matrix. This is consistent with the fact that $A$ is not diagonally dominant and that $\rho(B_J)>1$ and $\rho(B_{GS})>1$.
\paragraph{Codes}
\lstinputlisting[language=Matlab]{home1_mat/problem20_18.m}
% \lstinputlisting[language=Matlab]{home1_mat/Jacobi.m}
% \lstinputlisting[language=Matlab]{home1_mat/GS.m}
\newpage
\section{Problem20.19: Comparison of the results}
The numerical results are listed below (all logarithms are natural logarithms, i.e., base $e$):
\begin{itemize}
  \item \textbf{Jacobi:} converges in 30 iterations, with log-residual $-33.029204$.
  \item \textbf{Gauss--Seidel:} converges in 21 iterations, with log-residual $-33.886495$.
  \item \textbf{SOR$(\omega=1.1)$:} converges in 21 iterations, with log-residual $-33.365732$.
  \item \textbf{SOR$(\omega=1.2)$:} converges in 28 iterations, with log-residual $-33.997858$.
  \item \textbf{SOR$(\omega=1.3)$:} converges in 35 iterations, with log-residual $-33.057573$.
  \item \textbf{SOR$(\omega=1.5)$:} converges in 60 iterations, with log-residual $-33.292912$.
  \item \textbf{SOR$(\omega=1.9)$:} does not converge.
  \item \textbf{$B\backslash c$:} log-residual $-37.450359$.
\end{itemize}

From these results, we see that the Gauss--Seidel method converges faster than the Jacobi method, since it requires only 21 iterations instead of 30. For the SOR method, the choice of $\omega$ has a strong effect on performance. Among the tested values, $\omega=1.1$ works best in terms of convergence speed, since it converges in 21 iterations, the same as Gauss--Seidel. As $\omega$ increases, the convergence becomes slower, and when $\omega=1.9$, the method does not converge.

In terms of accuracy, the direct solver $B\backslash c$ gives the smallest residual, since its log-residual is $-37.450359$, which is smaller than those of all iterative methods. Therefore, $B\backslash c$ produces the most accurate solution overall.

In conclusion, among the iterative methods, Gauss--Seidel and SOR with $\omega=1.1$ perform best, while the MATLAB command $B\backslash c$ gives the highest accuracy.
\paragraph{Codes}
% \lstinputlisting[language=Matlab]{home1_mat/problem20_19.m}
\end{document}
